The baseline sweep has three main implications for ARC as an RL domain. First, the algorithm ordering is stable across task families. PPO outperforms DQN, PQN, and REINFORCE on every subset tested, and REINFORCE is a consistent runner-up. This ordering holds across the ARC-AGI and ConceptARC groups, despite the large structural differences between these task families. It suggests that the on-policy, advantage-estimated gradient signal matters more for ARC than the replay buffer and target-network machinery of value-based methods, at least under the architecture and hyperparameters used here.

Second, task-family variation dwarfs algorithm variation. The difference between PPO on Count (episodic return above 0.9) and PPO on Hard (episodic return near 0.1) is far larger than the difference between PPO and any other algorithm on a fixed subset. This means that aggregate metrics over the entire ARC corpus mask the structure of the problem. A useful set of baselines for ARC must report per-family or per-concept results, and JaxARC's throughput makes this level of granularity practical.

Third, the observation wrapper analysis on MiniARC shows that the agent interface matters. When the answer grid is visible, learnability improves. But simply concatenating more channels (input grids, contextual support examples) does not monotonically help. This is consistent with the hypothesis that the bottleneck is the network's capacity to extract and use relational information, not the amount of raw data in the observation.

\paragraph{Limitations.}
These baselines use a single, deliberately simple architecture: a shallow two-layer CNN with 32 channels. We chose this architecture to keep experiments tractable and to provide a clean lower bound, but stronger architectures could shift both the algorithm ranking and the absolute performance levels. The sweep also does not include meta-learning or test-time adaptation methods, which represent a separate and increasingly active line of work on ARC \citep{coleDontThrowBaby2025,akyurekSurprisingEffectivenessTestTime2025}. The current results characterize single-task training from scratch and should not be taken as upper bounds on what RL can achieve.

The four algorithms we tested all belong to model-free RL. Model-based methods, search-augmented agents, and hybrid neuro-symbolic approaches remain outside the scope of this sweep. Similarly, the ConceptARC and ARC-AGI subsets we report on are representative but not exhaustive; extended results covering additional subsets are in the appendix.

\paragraph{Future work.}
The baselines and the environment together define a clear experimental surface for future research. Stronger architectures (graph neural networks, transformers, or architectures with explicit compositional priors) can be evaluated against the PPO reference curves. Test-time adaptation and meta-learning methods can use JaxARC's throughput to run the large-scale sweeps these techniques require. The per-subset breakdowns reported here can guide task selection: Count and Entry are relatively easy targets for showing new methods work, while Hard and Tedious represent genuinely difficult ceilings. JaxARC removes the simulation bottleneck so that future work can focus on algorithms and representations rather than on engineering faster environments.
